\documentclass[10pt,a4paper]{amsart}
\usepackage[margin=19mm]{geometry}
\usepackage{amsmath,amssymb,mathtools}
\usepackage[scr=rsfs,cal=euler]{mathalfa}
\usepackage{xfrac}
\usepackage[unicode,hidelinks,bookmarks=false]{hyperref}
\newcommand{\ie}{i.e.\ }
\newcommand{\cf}{cf.\ }
\newcommand{\resp}{respectively\ }
\newcommand{\Subsection}{\subsection}
\providecommand{\nobreakdash}{\nobreak\hbox{-}}
\setlength{\emergencystretch}{2em}
\multlinegap0pt
\usepackage{tikz}
\usetikzlibrary{cd}
\usepackage{amssymb}
\usepackage{dsfont}
\usepackage{mathtools}
\usepackage[shortlabels]{enumitem}
\newtheorem{proposition}{Proposition}
\newcommand{\daq}[4]{\operatorname{D}_{#1}(#2|#3;#4)}
\newcommand{\n}{\mathfrak{n}}
\DeclareMathOperator{\uppercidim}{CI^{*}dim}
\title{Large Homomorphisms on the Homotopy Lie Coalgebra}
\author{Andrew J. Soto Levins, Ryan Watson}
\date{Source: arXiv:2605.03813v1 (CC BY 4.0)}
\begin{document}
\maketitle

\noindent\emph{Source and licence.} Large Homomorphisms on the Homotopy Lie Coalgebra. Version arXiv:2605.03813v1 (CC BY 4.0), distributed by arXiv under the Creative Commons Attribution 4.0 International licence, http://creativecommons.org/licenses/by/4.0/, as recorded in the arXiv licence metadata for this exact version and shown on the abstract page https://arxiv.org/abs/2605.03813v1. Full licence notice: \texttt{licenses/arxiv-2605.03813-CC-BY-4.0.txt}.\par\smallskip

\noindent\textit{Editorial scope.} Counted unit: the article's proposition P\_QuillenConjecture (source0.tex lines 484--485). The article's complete original proof of it is delivered with it (source0.tex lines 487--503, one \texttt{proof} environment of 17 lines). No mathematics is written, repaired or re-derived: the statement and the proof are the article's own lines, in the article's own notation, and no retained line is substituted or re-flowed.  No further source-local context is needed: every cross-reference of every retained line resolves inside the two delivered spans.  Nothing was omitted from any retained span, so the package carries no omission record.  The retained lines cite 5 works, whose entries are transcribed verbatim from the article's own bibliography source in the e-print and delivered with short editorial labels recorded in the item's reference mapping.  The retained lines contain the article's own diagram code, which the wrapper typesets with the standard tikz packages; no image file is delivered and the package declares no asset dependency.  Editorial formatting only, in addition: an AMS \texttt{amsart} preamble that restates the article's own declarations and macros used by the retained lines, namely the environments \texttt{proposition} on separate counters, together with the standard packages those lines need.  The retained environments print in this document as Proposition 1 in order of appearance, whereas the article prints the same items with the numbering of its own full document.  The complete native source of the article as distributed in the arXiv e-print for this exact version is bundled unchanged as \texttt{sources/199816/source0.tex}.
\begin{proposition} \label{P_QuillenConjecture} Let $\varphi\colon R\rightarrow (S,\n,\ell)$ be a local homomorphism with $\uppercidim{\varphi}$ finite. If $\daq{n}{S}{R}{\ell}=0$ for some $n\geq 2$, then $\varphi$ is q.c.i. at $\n$.
\end{proposition}

\begin{proof}
Let 
\[\begin{tikzcd}
	&& {R'} && {\widehat{S}} \\
	Q && {R''} && {\widehat{S}\otimes_{R'}R''}
	\arrow["{\varphi'}", from=1-3, to=1-5]
	\arrow[from=1-3, to=2-3]
	\arrow[from=1-5, to=2-5]
	\arrow["f", from=2-1, to=2-3]
	\arrow["g", from=2-3, to=2-5]
\end{tikzcd}\]
be a diagram witnessing the finiteness of the upper complete intersection dimension of $\varphi$ and let $L$ be the residue field of $\widehat{S}\otimes_{R'}R''$. Since $\ell$ is a subfield of $L$, $\daq{n}{S}{R}{L}$ is a direct sum of $\daq{n}{S}{R}{\ell}$, giving the equality below
\[\daq{n}{\widehat{S}\otimes_{R'}R''}{R''}{L} \cong \daq{n}{\widehat{S}}{R'}{L} \cong \daq{n}{S}{R}{L} = 0.\]
The first isomorphism is by \cite[6.3]{Iyengar:2007} and the second isomorphism is by \cite[Lemma 1.7]{Avramov:1999}. Since $f$ is c.i., $\daq{i}{R''}{Q}{L}=0$ for $i\geq 2$, and so from the Jacobi-Zariski exact sequence \cite[Theorem 5.1]{Andre:1974}
\[\daq{n}{R''}{Q}{L}\rightarrow \daq{n}{\widehat{S}\otimes_{R'}R''}{Q}{L}\rightarrow \daq{n}{\widehat{S}\otimes_{R'}R''}{R''}{L}\]
we get $\daq{n}{\widehat{S}\otimes_{R'}R''}{Q}{L} = 0$. The implies $gf$ is c.i. by \cite[Theorem A]{Briggs/Iyengar:2020} and \cite[Proposition 4.57]{Andre:1974}. Now by \cite[8.9]{Avramov/Henriques/Sega:2013} $g$ is q.c.i., which implies $\varphi$ is q.c.i. at $\n$ by \cite[8.7]{Avramov/Henriques/Sega:2013}.
\end{proof}

\begin{thebibliography}{99}
\bibitem[Andre:]{Andre:1974} Michel Andrè. \newblock {\em Homologie des algèbres commutatives}. \newblock Springer Berlin Heidelberg, 1974.
\bibitem[Avramo]{Avramov/Henriques/Sega:2013} Luchezar~L. Avramov, Inês~B.A. Henriques, and Liana~M. {\c S}ega. \newblock Quasi-complete intersection homomorphisms. \newblock {\em Pure and Applied Mathematics Quarterly}, 9(4):579--612, 2013.
\bibitem[Avramo]{Avramov:1999} Luchezar~L. Avramov. \newblock Locally complete intersection homomorphisms and a conjecture of {Q}uillen on the vanishing of cotangent homology. \newblock {\em Ann. of Math. (2)}, 150(2):455--487, 1999.
\bibitem[Briggs]{Briggs/Iyengar:2020} Benjamin Briggs and Srikanth Iyengar. \newblock Rigidity properties of the cotangent complex. \newblock {\em J. Amer. Math. Soc.}, 36(1):291--310, 2023.
\bibitem[Iyenga]{Iyengar:2007} Srikanth Iyengar. \newblock Andr\'{e}-{Q}uillen homology of commutative algebras. \newblock In {\em Interactions between homotopy theory and algebra}, volume 436 of {\em Contemp. Math.}, pages 203--234. Amer. Math. Soc., Providence, RI, 2007.
\end{thebibliography}
\end{document}
